% Generated draft: requires mathematical and source review.
\documentclass[11pt]{article}
\usepackage[margin=25mm]{geometry}
\usepackage{fontspec}
\setmainfont{Latin Modern Roman}
\usepackage{amsmath,amssymb,mathtools,unicode-math}
\setmathfont{Latin Modern Math}
\usepackage{enumitem,needspace,xurl,hyperref}
\setlength{\parindent}{0pt}
\setlength{\parskip}{5pt}
\setlength{\emergencystretch}{3em}
\newcommand{\N}{\mathbb{N}}
\newcommand{\Z}{\mathbb{Z}}
\newcommand{\Q}{\mathbb{Q}}
\newcommand{\R}{\mathbb{R}}
\newcommand{\C}{\mathbb{C}}
\newcommand{\F}{\mathbb{F}}
\newcommand{\A}{\mathbb{A}}
\newcommand{\PP}{\mathbb{P}}
\newcommand{\E}{\mathbb{E}}
\newcommand{\eps}{\varepsilon}
\newcommand{\dd}{\,\mathrm{d}}
\begin{document}
\section*{UnsolvedMath ID: 10400147}
\textbf{Source problem code:} catalogue entry
\subsection{Definitions and mathematical statement}
The source asks to find and classify all topological quantum field theories (TQFTs); the summary identifies the intended dimension as three, although the excerpt itself does not state it.

A standard nonextended oriented three-dimensional formulation, over a fixed field \(k\), is a symmetric monoidal functor
\[
 Z:\operatorname{Bord}^{\mathrm{or}}_3\longrightarrow\operatorname{Vect}^{\mathrm{fd}}_k.
\]
Objects are closed oriented surfaces, morphisms are oriented three-dimensional cobordisms up to the specified boundary-preserving diffeomorphism, composition is gluing and the monoidal operation is disjoint union. The functor assigns finite-dimensional state spaces and compatible linear maps, carrying disjoint union to tensor product and the empty surface to \(k\). In that formulation the classification target is all such functors up to symmetric monoidal natural isomorphism.

This is an explicit conventional reading, not an imposed full scope. Dimension, orientation/framing or other tangential structures, coefficient category, unitarity, anomalies/projective theories, and extended versus nonextended structure must be fixed before one classification problem is identified. No assumption that a chosen algebraic construction exhausts all theories is added. Finding examples and proving an exhaustive classification remain distinct requests.

The summary labels the unrestricted three-dimensional task a major open problem, but provides no precise classification theorem or list of admissible theories. That source assessment is retained without claiming that every restricted or extended version has the same status.
\subsection{Short English statement}
Find and classify the intended three-dimensional TQFTs. A usual functorial formulation is given, but the source omits the dimension/structure, target category and extendedness choices that determine the actual classification problem.
\subsection{Sources}
\begin{itemize}
\item Local registry: \url{data/unsolvedmath/problems.json}, permanent UnsolvedMath ID \(10400147\), source code catalogue entry; supplied statement, background, and research summary, updated 2026-08-06.
\item Ohtsuki, \emph{Problems on invariants of knots and 3-manifolds} (2002), Problem 8.1, PDF page 122: \url{https://msp.org/gtm/2002/04/gtm-2002-04-024s.pdf}. Supplied PDF-font-structure extraction accessed 2026-07-29; provenance distinguishes 231 numbered items.
\item No separate exact TQFT-scope definition or current classification reference is supplied.
\end{itemize}
\end{document}
